\chapter{Kernmodelle und Stabilitaet}

\section{Ziel dieses Kapitels}

Dieses Kapitel behandelt Modelle, mit denen man Atomkerne beschreibt.
Kapitel 3 hat die Bindungsenergie und die Weizsaecker-Massenformel behandelt.
Diese Formel erklaert globale Massentrends sehr gut, aber sie erklaert nicht alle
Details der Kernstabilitaet.

Insbesondere erklaert sie nicht vollstaendig, warum bestimmte Protonen- und
Neutronenzahlen besonders stabil sind.
Dafuer braucht man Kernmodelle.

\begin{notebox}
Dieses Kapitel behandelt die Stichwoerter:
\begin{itemize}
    \item Kernmodelle,
    \item Troepfchenmodell,
    \item Fermigasmodell,
    \item Potentialtopfmodell,
    \item Schalenmodell,
    \item magische Zahlen,
    \item Spin-Bahn-Kopplung,
    \item Separationsenergie,
    \item Paarungseffekte,
    \item Kerndeformation,
    \item Quadrupolmoment,
    \item kollektive Modelle,
    \item Stabilitaetsinsel.
\end{itemize}
\end{notebox}

\section{Warum braucht man Kernmodelle?}

Ein Atomkern ist ein quantenmechanisches Vielteilchensystem.
Er besteht aus Protonen und Neutronen, die ueber die starke Wechselwirkung
miteinander wechselwirken.

\begin{conceptbox}
Ein vollstaendig exaktes Modell eines schweren Kerns waere extrem kompliziert,
weil viele Nukleonen gleichzeitig miteinander wechselwirken.

Deshalb verwendet man Modelle.
Ein Modell beschreibt nicht jedes Detail exakt, sondern erklaert bestimmte
Eigenschaften besonders gut.
\end{conceptbox}

\begin{notebox}
Wichtige Fragen an Kernmodelle sind:
\begin{itemize}
    \item Warum sind Kerne gebunden?
    \item Warum gibt es stabile und instabile Kerne?
    \item Warum sind bestimmte Protonen- oder Neutronenzahlen besonders stabil?
    \item Warum besitzen Kerne diskrete Anregungszustaende?
    \item Warum sind manche Kerne kugelfoermig und andere deformiert?
\end{itemize}
\end{notebox}

\section{Globale und mikroskopische Modelle}

\begin{definitionbox}
Ein \emph{globales Kernmodell} beschreibt allgemeine Trends ueber viele Kerne hinweg.

Ein \emph{mikroskopisches Kernmodell} versucht staerker zu beschreiben, wie einzelne
Nukleonen quantenmechanische Zustaende besetzen.
\end{definitionbox}

\begin{conceptbox}
Das Troepfchenmodell beziehungsweise die Weizsaecker-Massenformel ist ein globales
Modell.
Es beschreibt mittlere Trends der Bindungsenergie.

Das Schalenmodell ist mikroskopischer.
Es erklaert, warum bestimmte Nukleonenzahlen besonders stabil sind.
\end{conceptbox}

\begin{examtipbox}
Typische Pruefungsfrage:

Warum reicht die Weizsaecker-Massenformel allein nicht aus?

Antwort:
Die Weizsaecker-Massenformel erklaert globale Trends der Bindungsenergie.
Sie erklaert aber nicht die besondere Stabilitaet bei magischen Zahlen und nicht die
detaillierte Struktur einzelner Kernzustaende.
\end{examtipbox}

\section{Das Troepfchenmodell}

Das Troepfchenmodell beschreibt den Kern aehnlich wie einen Tropfen einer Fluessigkeit.
Es ist die anschauliche Grundlage der Weizsaecker-Massenformel.

\begin{definitionbox}
Das \emph{Troepfchenmodell} behandelt den Atomkern als annahernd inkompressiblen
Tropfen aus Kernmaterie.
\end{definitionbox}

\begin{conceptbox}
Die Analogie zu einem Fluessigkeitstropfen erklaert mehrere Beitraege zur
Bindungsenergie:
\begin{itemize}
    \item Volumenenergie: Nukleonen im Inneren sind gebunden.
    \item Oberflaechenenergie: Nukleonen an der Oberflaeche haben weniger Nachbarn.
    \item Coulombenergie: Protonen stossen sich elektrisch ab.
    \item Asymmetrieenergie: ein starkes Ungleichgewicht zwischen $N$ und $Z$ kostet Energie.
    \item Paarungsenergie: gepaarte Nukleonen sind energetisch guenstiger.
\end{itemize}
\end{conceptbox}

\begin{formulabox}
Eine uebliche Form der Weizsaecker-Bindungsenergie ist:
\begin{equation*}
    B(A,Z)
    =
    a_V A
    - a_S A^{2/3}
    - a_C \frac{Z(Z-1)}{A^{1/3}}
    - a_A \frac{(N-Z)^2}{A}
    + \delta(A,Z).
\end{equation*}
\end{formulabox}

\begin{notebox}
Das Troepfchenmodell erklaert gut:
\begin{itemize}
    \item den groben Verlauf der Bindungsenergie pro Nukleon,
    \item die ungefaehre Lage stabiler Kerne,
    \item Massenparabeln bei Isobaren,
    \item qualitative Aspekte der Kernspaltung.
\end{itemize}

Es erklaert aber nicht die magischen Zahlen.
\end{notebox}

\section{Das Fermigasmodell}

Das Fermigasmodell betrachtet Protonen und Neutronen als Fermionen, die Zustaende
in einem effektiven Kernvolumen besetzen.

\begin{definitionbox}
Im \emph{Fermigasmodell} werden die Nukleonen als Fermionen in einem gemeinsamen
Kernvolumen behandelt.

Protonen und Neutronen fuellen quantenmechanische Zustaende bis zu einer
Fermi-Energie.
\end{definitionbox}

\begin{conceptbox}
Protonen und Neutronen haben Spin
\begin{equation*}
    s=\frac{1}{2}.
\end{equation*}
Sie sind daher Fermionen und unterliegen dem Pauli-Prinzip.

Das Pauli-Prinzip besagt:
Zwei identische Fermionen koennen nicht denselben quantenmechanischen Zustand besetzen.
\end{conceptbox}

\begin{definitionbox}
Die \emph{Fermi-Energie} $E_F$ ist die Energie des hoechsten besetzten Zustands
bei Temperatur nahe null.
\end{definitionbox}

\begin{conceptbox}
Im Kern gibt es getrennte Fermi-Niveaus fuer Protonen und Neutronen.

Wenn sehr viel mehr Neutronen als Protonen vorhanden sind, muessen die Neutronen
hoehere Zustaende besetzen.
Das erhoeht die kinetische Energie.

Diese Idee fuehrt zum Asymmetrieterm der Weizsaecker-Massenformel:
\begin{equation*}
    -a_A\frac{(N-Z)^2}{A}.
\end{equation*}
\end{conceptbox}

\begin{examtipbox}
Das Fermigasmodell liefert eine anschauliche Begruendung fuer den Asymmetrieterm:

Wenn $N\neq Z$ gilt, sind die Fermi-Niveaus von Neutronen und Protonen ungleich
gefuellt.
Das erhoeht die kinetische Energie und verringert dadurch die Bindung.
\end{examtipbox}

\section{Potentialtopfmodell}

Im Potentialtopfmodell bewegt sich jedes Nukleon in einem mittleren Potential,
das durch alle anderen Nukleonen erzeugt wird.

\begin{definitionbox}
Das \emph{Potentialtopfmodell} beschreibt ein Nukleon so, als bewege es sich in
einem effektiven Potentialtopf des Kerns.
\end{definitionbox}

\begin{conceptbox}
Obwohl Nukleonen miteinander wechselwirken, kann man naeherungsweise annehmen:
Jedes einzelne Nukleon bewegt sich in einem mittleren Potential, das von allen
anderen Nukleonen erzeugt wird.

Diese Naeherung fuehrt auf diskrete Energieniveaus.
\end{conceptbox}

\begin{notebox}
In den Vorlesungsabbildungen werden Protonen und Neutronen durch getrennte
Potentialtoepfe dargestellt.

Fuer Protonen ist das Potential wegen der Coulomb-Abstossung veraendert.
Deshalb unterscheiden sich Protonen- und Neutronenpotentiale.
\end{notebox}

\begin{conceptbox}
Fuer Neutronen wirkt im Wesentlichen das Kernpotential.

Fuer Protonen wirkt zusaetzlich das Coulomb-Potential.
Dadurch liegt das Protonenpotential effektiv hoeher beziehungsweise wird am Rand
anders geformt.
\end{conceptbox}

\section{Stehende Wellen im Kasten}

Ein einfaches quantenmechanisches Bild fuer gebundene Nukleonen ist das Teilchen
im Kasten.

\begin{conceptbox}
Ein gebundenes Teilchen kann nicht beliebige Energien besitzen.
Es bildet stehende Wellen im Potential.

Nur Wellenfunktionen, die zu den Randbedingungen passen, sind erlaubt.
Daraus entstehen diskrete Energieniveaus.
\end{conceptbox}

\begin{formulabox}
Fuer ein einfaches Teilchen im eindimensionalen Kasten der Laenge $L$ sind die
erlaubten Wellenzahlen:
\begin{equation*}
    k_n = \frac{n\pi}{L},
    \qquad
    n=1,2,3,\ldots
\end{equation*}
Die Energien sind proportional zu
\begin{equation*}
    E_n \propto k_n^2.
\end{equation*}
\end{formulabox}

\begin{conceptbox}
Im dreidimensionalen Kasten gibt es Zustaende mit Quantenzahlen
\begin{equation*}
    n_x,\qquad n_y,\qquad n_z.
\end{equation*}

Die Zustaende koennen als Gitterpunkte im Impuls- beziehungsweise $k$-Raum gezaehlt
werden.
Dieses Bild ist wichtig fuer das Fermigasmodell.
\end{conceptbox}

\begin{examtipbox}
Die zentrale Aussage ist:

Raeumliche Begrenzung fuehrt zu stehenden Wellen.
Stehende Wellen fuehren zu diskreten Energieniveaus.
\end{examtipbox}

\section{Vom Potentialtopf zum Schalenmodell}

Das Schalenmodell entsteht aus der Idee, dass Nukleonen diskrete Einteilchenzustaende
in einem mittleren Kernpotential besetzen.

\begin{definitionbox}
Das \emph{Schalenmodell} beschreibt Atomkerne durch die Besetzung diskreter
Energiezustaende durch Protonen und Neutronen.
\end{definitionbox}

\begin{conceptbox}
Das Schalenmodell ist analog zur Elektronenhuelle in der Atomphysik.

In der Atomphysik erklaeren abgeschlossene Elektronenschalen die besondere Stabilitaet
der Edelgase.

In der Kernphysik erklaeren abgeschlossene Protonen- oder Neutronenschalen die
besondere Stabilitaet bestimmter Kerne.
\end{conceptbox}

\begin{conceptbox}
Protonen und Neutronen besetzen getrennte Schalen.

Eine abgeschlossene Protonenschale bedeutet:
\begin{equation*}
    Z = \text{magische Zahl}.
\end{equation*}

Eine abgeschlossene Neutronenschale bedeutet:
\begin{equation*}
    N = \text{magische Zahl}.
\end{equation*}
\end{conceptbox}

\section{Magische Zahlen}

Bestimmte Protonen- und Neutronenzahlen fuehren zu besonders stabilen Kernen.
Diese Zahlen heissen magische Zahlen.

\begin{definitionbox}
\emph{Magische Zahlen} sind Protonen- oder Neutronenzahlen, bei denen eine Kernschale
abgeschlossen ist.

Kerne mit magischer Protonen- oder Neutronenzahl sind besonders stabil.
\end{definitionbox}

\begin{formulabox}
Die wichtigsten magischen Zahlen sind:
\begin{equation*}
    2,\quad 8,\quad 20,\quad 28,\quad 50,\quad 82,\quad 126.
\end{equation*}
\end{formulabox}

\begin{conceptbox}
Wenn $Z$ magisch ist, ist die Protonenschale abgeschlossen.

Wenn $N$ magisch ist, ist die Neutronenschale abgeschlossen.

Wenn sowohl $Z$ als auch $N$ magisch sind, spricht man von einem doppelt magischen
Kern.
\end{conceptbox}

\begin{examplebox}
Beispiele fuer doppelt magische Kerne sind:
\begin{equation*}
    {}^{4}_{2}\mathrm{He}
    \qquad
    (Z=2,\;N=2),
\end{equation*}
\begin{equation*}
    {}^{16}_{8}\mathrm{O}
    \qquad
    (Z=8,\;N=8),
\end{equation*}
\begin{equation*}
    {}^{40}_{20}\mathrm{Ca}
    \qquad
    (Z=20,\;N=20),
\end{equation*}
\begin{equation*}
    {}^{208}_{82}\mathrm{Pb}
    \qquad
    (Z=82,\;N=126).
\end{equation*}
\end{examplebox}

\begin{examtipbox}
Magische Zahlen sind Pruefungswissen.
Man sollte mindestens die Liste kennen:
\begin{equation*}
    2,\;8,\;20,\;28,\;50,\;82,\;126.
\end{equation*}

Ausserdem sollte man erklaeren koennen:
Eine magische Zahl bedeutet eine abgeschlossene Schale.
Eine abgeschlossene Schale bedeutet besondere Stabilitaet.
\end{examtipbox}

\section{Warum einfache Potentiale nicht ausreichen}

Ein einfacher dreidimensionaler harmonischer Oszillator liefert zwar Schalen,
aber nicht die experimentell richtigen magischen Zahlen.

\begin{conceptbox}
Einfache Potentiale wie der harmonische Oszillator erklaeren einen Teil der
Schalenstruktur.
Sie liefern aber nicht alle beobachteten magischen Zahlen.

Insbesondere die magischen Zahlen
\begin{equation*}
    28,\quad 50,\quad 82,\quad 126
\end{equation*}
benoetigen eine zusaetzliche wichtige Wechselwirkung:
die Spin-Bahn-Kopplung.
\end{conceptbox}

\begin{notebox}
Die Vorlesungsabbildungen zu Oszillatorniveaus zeigen Schalenstrukturen.
Fuer die tatsaechlich beobachteten magischen Zahlen muss das Modell aber durch
Spin-Bahn-Kopplung verbessert werden.
\end{notebox}

\section{Spin-Bahn-Kopplung}

Die Spin-Bahn-Kopplung koppelt den Spin eines Nukleons an seinen Bahndrehimpuls.

\begin{definitionbox}
Die \emph{Spin-Bahn-Kopplung} ist eine Wechselwirkung zwischen dem Spin $\vec{s}$
eines Nukleons und seinem Bahndrehimpuls $\vec{l}$.

Sie ist proportional zu
\begin{equation*}
    \vec{l}\cdot\vec{s}.
\end{equation*}
\end{definitionbox}

\begin{conceptbox}
Der Gesamtdrehimpuls eines einzelnen Nukleons ist
\begin{equation*}
    \vec{j}=\vec{l}+\vec{s}.
\end{equation*}

Da Nukleonen Spin
\begin{equation*}
    s=\frac{1}{2}
\end{equation*}
haben, kann fuer gegebenes $l$ gelten:
\begin{equation*}
    j=l+\frac{1}{2}
\end{equation*}
oder
\begin{equation*}
    j=l-\frac{1}{2}.
\end{equation*}
\end{conceptbox}

\begin{conceptbox}
Die Spin-Bahn-Kopplung spaltet Energieniveaus auf.

Diese Aufspaltung verschiebt die Schalenabschluesse so, dass die experimentellen
magischen Zahlen erklaert werden koennen.
\end{conceptbox}

\begin{examtipbox}
Ohne Spin-Bahn-Kopplung erhaelt man nicht die richtige Folge der magischen Zahlen.

Die Spin-Bahn-Kopplung ist daher ein zentrales Element des Kernschalenmodells.
\end{examtipbox}

\section{Schalenabschluesse und Stabilitaet}

Ein Schalenabschluss bedeutet, dass eine Gruppe von Zustaenden vollstaendig besetzt ist.
Danach kommt eine Energieluecke zum naechsten Zustand.

\begin{conceptbox}
Wenn eine Schale abgeschlossen ist, ist der Kern besonders stabil.

Der Grund ist:
Um ein weiteres Nukleon hinzuzufuegen oder ein Nukleon anzuregen, muss man eine
groessere Energieluecke ueberwinden.
\end{conceptbox}

\begin{definitionbox}
Die \emph{Separationsenergie} ist die Energie, die benoetigt wird, um ein Nukleon
aus dem Kern zu entfernen.
\end{definitionbox}

\begin{conceptbox}
Bei magischen Zahlen treten oft besonders grosse Separationsenergien auf.
Das ist ein experimentelles Signal fuer Schalenabschluesse.
\end{conceptbox}

\begin{examtipbox}
Wenn eine Aufgabe fragt, woran man magische Zahlen erkennt, sind typische Hinweise:
\begin{itemize}
    \item hohe Stabilitaet,
    \item grosse Separationsenergie,
    \item viele stabile Isotope oder Isotone,
    \item niedrige Anregbarkeit,
    \item abgeschlossene Schale.
\end{itemize}
\end{examtipbox}

\section{Gerade-gerade, gerade-ungerade und ungerade-ungerade Kerne}

Neben Schalenabschluessen spielen Paarungseffekte eine wichtige Rolle.

\begin{definitionbox}
Ein \emph{gerade-gerade Kern} hat gerade Protonenzahl $Z$ und gerade Neutronenzahl $N$.

Ein \emph{ungerade-ungerade Kern} hat ungerade Protonenzahl $Z$ und ungerade
Neutronenzahl $N$.
\end{definitionbox}

\begin{conceptbox}
Nukleonen bilden bevorzugt Paare mit entgegengesetztem Spin.

Daher sind gerade-gerade Kerne besonders stabil.
Ungerade-ungerade Kerne sind meistens weniger stabil.
\end{conceptbox}

\begin{notebox}
Dieser Effekt steckt in der Weizsaecker-Massenformel im Paarungsterm
\begin{equation*}
    \delta(A,Z).
\end{equation*}
\end{notebox}

\begin{examtipbox}
Gerade-gerade Kerne sind besonders stabil.

Ungerade-ungerade Kerne sind meist weniger stabil.
\end{examtipbox}

\section{Kernformen und Deformation}

Nicht alle Kerne sind kugelfoermig.
Manche Kerne sind deformiert.

\begin{definitionbox}
Eine \emph{Kerndeformation} liegt vor, wenn die Form des Kerns von einer Kugel
abweicht.
\end{definitionbox}

\begin{conceptbox}
Kerne in der Naehe magischer Zahlen sind haeufig eher kugelfoermig, weil
abgeschlossene Schalen besonders symmetrisch und stabil sind.

Kerne weit weg von Schalenabschluessen koennen deformiert sein.
Typische Formen sind:
\begin{itemize}
    \item prolat, also laenglich,
    \item oblat, also abgeflacht.
\end{itemize}
\end{conceptbox}

\begin{notebox}
In den Vorlesungsabbildungen zu Kernmodellen treten Deformationsparameter und
Quadrupolmomente auf.
Diese Groessen beschreiben, wie stark ein Kern von der Kugelform abweicht.
\end{notebox}

\section{Quadrupolmoment}

Ein wichtiges Mass fuer Kernverformung ist das Quadrupolmoment.

\begin{definitionbox}
Das \emph{Quadrupolmoment} beschreibt die Abweichung einer Ladungsverteilung von
Kugelsymmetrie.
\end{definitionbox}

\begin{conceptbox}
Ein kugelsymmetrischer Kern hat kein intrinsisches Quadrupolmoment.

Ein deformierter Kern besitzt ein Quadrupolmoment.
Das Vorzeichen und die Groesse geben Hinweise auf die Form des Kerns.
\end{conceptbox}

\begin{notebox}
Eine genaue Berechnung von Quadrupolmomenten ist modellabhaengig.
Fuer dieses Kapitel ist wichtig:
Quadrupolmomente sind experimentelle Hinweise auf Deformation.
\end{notebox}

\section{Kollektive Modelle}

Das Schalenmodell beschreibt einzelne Nukleonen in einem mittleren Potential.
Manche Eigenschaften von Kernen sind aber kollektive Bewegungen vieler Nukleonen.

\begin{definitionbox}
Ein \emph{kollektives Modell} beschreibt Bewegungen, an denen viele Nukleonen
gemeinsam beteiligt sind.
\end{definitionbox}

\begin{conceptbox}
Wichtige kollektive Bewegungen sind:
\begin{itemize}
    \item Rotation deformierter Kerne,
    \item Vibrationen der Kernoberflaeche,
    \item kollektive Deformationen.
\end{itemize}
\end{conceptbox}

\begin{notebox}
Kollektive Modelle sind besonders wichtig fuer deformierte Kerne.
Das Schalenmodell ist besonders stark bei der Erklaerung magischer Zahlen und
Einteilchenzustaende.
\end{notebox}

\section{Zusammenhang der Modelle}

Kein einzelnes Modell erklaert alle Eigenschaften aller Kerne perfekt.
Die Modelle ergaenzen sich.

\begin{conceptbox}
Die wichtigsten Modelle haben unterschiedliche Staerken:
\begin{itemize}
    \item Das Troepfchenmodell erklaert globale Bindungsenergien und Massen.
    \item Das Fermigasmodell erklaert den Asymmetrieterm und Fermi-Energie-Ideen.
    \item Das Potentialtopfmodell erklaert diskrete Einteilchenzustaende.
    \item Das Schalenmodell erklaert magische Zahlen.
    \item Kollektive Modelle erklaeren Rotationen, Vibrationen und Deformationen.
\end{itemize}
\end{conceptbox}

\begin{examtipbox}
Wenn gefragt wird, welches Modell wofuer geeignet ist:
\begin{itemize}
    \item Bindungsenergie global: Troepfchenmodell und Weizsaecker-Massenformel.
    \item Magische Zahlen: Schalenmodell mit Spin-Bahn-Kopplung.
    \item Asymmetrieterm: Fermigasmodell.
    \item Deformation und Rotation: kollektive Modelle.
\end{itemize}
\end{examtipbox}

\section{Stabilitaet im Massental}

Das Massental beschreibt, welche Kerne energetisch bevorzugt sind.

\begin{definitionbox}
Das \emph{Massental} ist die Darstellung der Kernmassen beziehungsweise
Bindungsenergien als Funktion von Protonenzahl $Z$ und Neutronenzahl $N$.

Stabile Kerne liegen im Tal der niedrigsten Massen beziehungsweise hoechsten Bindungen.
\end{definitionbox}

\begin{conceptbox}
Die Weizsaecker-Massenformel beschreibt die grobe Form des Massentals.

Das Schalenmodell erklaert zusaetzliche lokale Stabilitaeten bei magischen Zahlen.
\end{conceptbox}

\begin{notebox}
Kernstabilitaet ergibt sich daher aus mehreren Beitraegen:
\begin{itemize}
    \item globale Massentrends,
    \item Coulombenergie,
    \item Asymmetrieenergie,
    \item Paarung,
    \item Schalenabschluesse,
    \item Deformation.
\end{itemize}
\end{notebox}

\section{Stabilitaetsinsel}

Bei sehr schweren Kernen erwartet man, dass Schalenabschluesse trotz grosser
Coulomb-Abstossung zusaetzliche Stabilitaet erzeugen koennen.

\begin{definitionbox}
Die \emph{Stabilitaetsinsel} ist ein theoretisch erwarteter Bereich superschwerer
Kerne, in dem Schalenabschluesse zu vergleichsweise laengeren Lebensdauern fuehren.
\end{definitionbox}

\begin{conceptbox}
Sehr schwere Kerne sind wegen der starken Coulomb-Abstossung normalerweise instabil.

Wenn jedoch Protonen- und Neutronenzahlen in die Naehe magischer Zahlen kommen,
koennen Schalenabschluesse zusaetzliche Stabilitaet liefern.

Dadurch entsteht die Idee einer Stabilitaetsinsel.
\end{conceptbox}

\begin{notebox}
Die Stabilitaetsinsel zeigt, warum das Troepfchenmodell allein nicht ausreicht.
Ohne Schalenmodell wuerde man die zusaetzliche Stabilisierung durch Schalenabschluesse
nicht verstehen.
\end{notebox}

\section{Mini-Zusammenfassung}

\begin{notebox}
Die wichtigsten Punkte dieses Kapitels sind:
\begin{itemize}
    \item Kernmodelle erklaeren unterschiedliche Aspekte des Atomkerns.
    \item Das Troepfchenmodell beschreibt globale Bindungsenergien.
    \item Das Fermigasmodell erklaert die Rolle des Pauli-Prinzips und den Asymmetrieterm.
    \item Das Potentialtopfmodell fuehrt auf diskrete Einteilchenzustaende.
    \item Das Schalenmodell erklaert magische Zahlen.
    \item Magische Zahlen sind:
    \begin{equation*}
        2,\quad 8,\quad 20,\quad 28,\quad 50,\quad 82,\quad 126.
    \end{equation*}
    \item Spin-Bahn-Kopplung ist noetig, um die richtigen magischen Zahlen zu erhalten.
    \item Gerade-gerade Kerne sind besonders stabil.
    \item Kerne koennen kugelfoermig oder deformiert sein.
    \item Quadrupolmomente geben Hinweise auf Deformation.
    \item Kollektive Modelle beschreiben gemeinsame Bewegungen vieler Nukleonen.
    \item Kernstabilitaet ergibt sich aus globalen Massentrends plus Schalen- und Paarungseffekten.
\end{itemize}
\end{notebox}

\section{Pruefungscheck}

\begin{examtipbox}
Nach diesem Kapitel solltest du folgende Fragen beantworten koennen:
\begin{enumerate}
    \item Warum braucht man verschiedene Kernmodelle?
    \item Was beschreibt das Troepfchenmodell?
    \item Was erklaert das Fermigasmodell?
    \item Was ist die Fermi-Energie?
    \item Warum fuehrt $N\neq Z$ zu einem Energiebeitrag?
    \item Was ist das Potentialtopfmodell?
    \item Warum entstehen diskrete Energieniveaus?
    \item Was ist das Schalenmodell?
    \item Was sind magische Zahlen?
    \item Welche magischen Zahlen sollte man kennen?
    \item Was ist ein doppelt magischer Kern?
    \item Warum braucht man Spin-Bahn-Kopplung?
    \item Warum sind Schalenabschluesse besonders stabil?
    \item Was ist Separationsenergie?
    \item Warum sind gerade-gerade Kerne besonders stabil?
    \item Was bedeutet Kerndeformation?
    \item Was beschreibt ein Quadrupolmoment?
    \item Was sind kollektive Modelle?
    \item Wie haengen Massental und Schalenmodell zusammen?
    \item Was ist die Stabilitaetsinsel?
\end{enumerate}
\end{examtipbox}